
\newif\iftsenvAppendixEnvironmentStandalone
\ifdefined\tsenvDocsCombinedAppendixMode
\tsenvAppendixEnvironmentStandalonefalse
\def\tsenvAppendixEnvironmentStart{%
  \docTitle{Appendix: Environment Catalog}%
  \phantomsection\label{docstart:appendix_environment}%
  \section*{Environment Catalog}%
}
\def\tsenvAppendixEnvironmentEnd{}
\else\ifdefined\tsenvCombinedAppendixMode
\tsenvAppendixEnvironmentStandalonefalse
\def\tsenvAppendixEnvironmentStart{%
  \section{Environment Catalog}%
  \label{app:environment-catalog}%
}
\def\tsenvAppendixEnvironmentEnd{}
\else
\tsenvAppendixEnvironmentStandalonetrue
\documentclass[11pt]{article}
\usepackage{manual_docs_style}
\usepackage{graphicx}
\usepackage{placeins}
%\externaldocument{build/generate_question_stage}[generate_question_stage.pdf]

\def\tsenvAppendixEnvironmentStart{%
  \begin{document}%
  \docTitle{Appendix: Environment Catalog}
  \phantomsection\label{docstart:appendix_environment}%
  \section*{Environment Catalog}%
}
\def\tsenvAppendixEnvironmentEnd{\end{document}}
\fi\fi

\ifdefined\tsenvCombinedAppendixMode
  \newcommand{\interventionTerm}{intervention}
\else
  \newcommand{\interventionTerm}{\code{first_time_diff}}
\fi


\newenvironment{ParameterSrdTable}{
  \begingroup
  \footnotesize
  \renewcommand{\arraystretch}{1.12}
  \setlength{\tabcolsep}{5pt}
  \begin{longtable}{|P{0.42\textwidth}|P{0.20\textwidth}|P{0.22\textwidth}|}
  \caption{Parameter minimum threhsold distance}\\
  \hline
  \textbf{Parameter} & \textbf{\termMinSRDDistance} & \textbf{\termMinAbsDist} \\
  \hline
  \endfirsthead
  \hline
  \textbf{Parameter} & \textbf{\termMinSRDDistance} & \textbf{\termMinAbsDist} \\
  \hline
  \endhead
}{
  \hline
  \end{longtable}
  \endgroup
}


\tsenvAppendixEnvironmentStart
\iftsenvAppendixEnvironmentStandalone
These documents are matched in the codebase for each model \code{models/simulink/<MODEL>}
\begin{itemize}
  \item The low and high textual-description level entries match \code{description_levels.json::textual_description::low} and  \code{description_levels.json::textual_description::high}, except for the sentence `These time series were generated by', because it is added during \hyperref[sec:question-text-schema]{question generation}.
  \item `Initial states and parameters' table matches \code{description_levels.json::internal_naming_to_agent_facing_parameter} and \code{experiment_config.json::exposed_variables::parameters::*::allowed_interval} for the parameters and \code{experiment_config.json::exposed_variables::initial_states::*::allowed_interval} for the intial states.
  \item `Minimum intervention distance' table matches \code{experiment_config.json::parameters::*::}\termMinAbsDist and \code{experiment_config.json::parameters::*::}\termMinSRDDistance
  \item `Simulation Properties' table matches \code{experiment_config.json::parameters}
  \item \code{experiment_config.json}: 
  \begin{itemize}
    \item `Initial states and parameters' table matches \code{internal_naming_to_agent_facing_parameter} and \code{exposed_variables::parameters::*::allowed_interval}
  \end{itemize}
  \item \code{description_levels.json}:
  \begin{itemize}
    \item The low and high entries to \code{textual_description::low} and  \code{textual_description::high}, except for the sentence `These time series were generated by', because it is added during \hyperref[sec:question-text-schema]{question generation}.
  \end{itemize} 
\end{itemize}


\paragraph{Utility scripts}
If you run 
\begin{DocCode}
    python3 workflows/utils_scripts/update_enviroment_scripts.py <MODEL>
\end{DocCode}
You can overwrite the entries of \code{description_levels.json} and \code{experiment_config.json} with values from this document.

\fi

% Use \iftsenvAppendixEnvironmentStandalone ... \fi for content that should
% appear only in the standalone environment appendix, not in the paper appendix.
\par\bigskip
{\LARGE\bfseries\texttt{BallDrop}\par}
\phantomsection\label{app:environment:balldrop}
\medskip
{\large\bfseries Textual-description levels\par}
\textbf{Low textual context}
\begin{DocCode}
These time series were generated by simulating a physical process.
\end{DocCode}

\textbf{High textual context}
\begin{DocCode}
These time series were generated by simulating a 1D vertical point-mass ball model under gravity.
While the ball is above the ground, its motion is governed by gravity and quadratic air drag acting opposite the direction of travel, and the position evolves according to the current velocity.
Ground interaction is modeled with a restitution-based hard-stop law.
When the ball reaches the ground with impact speed equal or above a prescribed impact-speed cutoff 0.5 m/s, the impact is treated as instantaneous and the post-impact speed is set by the coefficient of restitution e in [0,1].
The rebound velocity points upward and its magnitude equals e times the pre-impact speed.
When the impact speed is below the threshold, the ball does not rebound and instead enters static contact at the ground.

\end{DocCode}

{\large\bfseries Relevant Tables\par}
\begin{EnvironmentConfigTableNoMod}
\textemdash & \code{initial_height} & \code{[0, 6]} \\
\textemdash & \code{initial_velocity} & \code{[-10, 2]} \\
\hline
mass & \code{mass} & \code{[0.5, 20]} \\
drag coefficient & \code{drag_coeff} & \code{[0, 30]} \\
coefficient of restitution & \code{restitution} & \code{[0.15, 1.0]} \\
gravity acceleration & \code{gravity} & \code{[2, 15]} \\
\end{EnvironmentConfigTableNoMod}

\begin{ParameterSrdTable}
\code{mass} & \code{0.1} & \code{1} \\
\code{drag_coeff} & \code{0.0} & \code{2} \\
\code{restitution} & \code{0.0} & \code{0.3} \\
\code{gravity} & \code{0.3} & \code{3} \\
\end{ParameterSrdTable}

\begin{SimulationPropertiesTable}{Simulation Properties}
\code{end_time_input_s} & \code{10.0} \\
\code{sampling_rate_hz} & \code{400.0} \\
\end{SimulationPropertiesTable}

\begin{MappingTable}{Observed Signals.}
col1: ball height & \code{Position}\\
col2: ball velocity & \code{Velocity} \\
col3: contact impulse (N*s) & \code{Hard_Stop_f} \\
col4: time & \code{time} \\
\end{MappingTable}

\paragraph{Problem-specific features.}
\begin{itemize}
  \item \code{bouncing_events}: List of impact times when the ball bounces on the ground.
  \item \code{apex_heights}: List of all the apex heights after each impact
  \item \code{is_terminal_velocity}: Time intervals during which the ball is approximately at terminal velocity.
  \item \code{max_velocity}: Max velocity of the ball
  \item \code{acceleration}: Derivative of the velocity
\end{itemize}

\paragraph{Environment-specific detectability rules}
\begin{itemize}
  \item After \interventionTerm, at least one of the baseline or intervention traces must have an apex with height greater than \code{0.1}.
  \item Every trace with usable impact data must have fewer than \code{25} detected rebounding bounces.
  \item If the correct class is "drag coefficient" then either:
  \begin{itemize}
    \item the normalized DTW between the estimated free-flight acceleration traces over \code{[first_time_diff, first_time_diff + 0.2]}, excluding samples near ground-contact events, must be larger than \code{0.1}.
    \item \code{is_terminal_velocity} is True after at least one bounce 
    \item \code{is_terminal_velocity} turns from True to False while the position is higher than \code{0.1}
  \end{itemize}
  Additional after \interventionTerm there must be at least one bounce.
  \item If the parameter change is "restitution", "mass", "drag coefficient" then the \interventionTerm time must be after a first bounce and one bounce after.
\end{itemize}

\paragraph{Signal preprocessing}
\begin{itemize}
  \item \code{Hard_Stop_f}:  We keep only peaks that look like a real ground-impact rebound. Specifically
  \begin{itemize}
      \item median Velocity in \code{[peak_time - 0.1s, peak_time]} must be \code{< 0}
      \item median Velocity in \code{[peak_time, peak_time + 0.1s]} must be \code{> 0}
  \end{itemize}
\end{itemize}

{\large\bfseries Plot\par}

\begin{figure}[ht]
\centering
\IfFileExists{plots/environments/9860761527be368b6e6e5315b3fb7a8d__gravity__Position__noise_stack_low_high.pdf}{%
  \includegraphics[width=0.85\linewidth]{plots/environments/9860761527be368b6e6e5315b3fb7a8d__gravity__Position__noise_stack_low_high.pdf}%
}{%
  \includegraphics[width=0.85\linewidth]{docs/plots/environments/9860761527be368b6e6e5315b3fb7a8d__gravity__Position__noise_stack_low_high.pdf}%
}
\caption{Position Channel from BallDrop.}
\label{fig:balldrop_env}
\end{figure}

\newpage

\par\bigskip
{\LARGE\bfseries\texttt{DampedMassBetweenWalls}\par}
\medskip
{\large\bfseries Textual-description levels\par}
\textbf{Low textual context}
\begin{DocCode}
These time series were generated by simulating a physical process.
\end{DocCode}

\textbf{High textual context}
\begin{DocCode}
These time series were generated by simulating a damped mass moving along a one-dimensional rail between two rigid walls located at positions x_L and x_R.
The rail may be tilted by a fixed inclination angle. 
When the mass is not in contact with either wall, its motion along the rail is governed by viscous drag, proportional to velocity and opposite the direction of motion, and, when the inclination angle is nonzero, by the component of gravity along the rail.
When the mass hits a wall, the collision is modeled as an instantaneous impact: the direction of motion reverses and the post-impact speed is reduced according to that walls' restitution coefficient. 

\end{DocCode}

\begin{EnvironmentConfigTableNoMod}
\textemdash & \code{initial_position} & \code{[3, 10]} \\
\textemdash & \code{initial_velocity} & \code{[-10, 10]} \\
\hline
weight of the mass & \code{mass} & \code{[0.1, 100]} \\
coefficient of viscous damping & \code{viscous_damping} & \code{[0.001, 12]} \\
coefficient of restitution of the right wall & \code{restitution_right} & \code{[0.3, 0.9999]} \\
coefficient of restitution of the left wall & \code{restitution_left} & \code{[0.25, 0.9999]} \\
inclination angle of the rail & \code{inclination_angle} & \code{[-4, 4]} \\
\end{EnvironmentConfigTableNoMod}

\begin{ParameterSrdTable}
\code{mass} & \code{0.3} & \code{0.06} \\
\code{viscous_damping} & \code{0.3} & \code{0.06} \\
\code{restitution_right} & \code{0.3} & \code{0.06} \\
\code{restitution_left} & \code{0.3} & \code{0.06} \\
\code{inclination_angle} & \code{0.3} & \code{0.06} \\
\end{ParameterSrdTable}

\begin{SimulationPropertiesTable}{Simulation Properties}
\code{end_time_input_s} & \code{60.0} \\
\code{sampling_rate_hz} & \code{50.0} \\
\end{SimulationPropertiesTable}

\begin{MappingTable}{Observed Signals. Sampling frequency: \code{50.0} Hz.}
col1: mass position along the rail & \code{Initial_Spacer_F_x} \\
col2: mass velocity along the rail & \code{Initial_Spacer_F_v} \\
col3: force on the right wall & \code{Right_Hard_Stop_f} \\
col4: time & \code{time} \\
\end{MappingTable}

\paragraph{Problem-specific features.}
\begin{itemize}
  \item \code{left_impact_times}: List of estimated impact times at the left wall.
  \item \code{right_impact_times}: List of estimated impact times at the right wall.
  \item \code{restitution_left_estimates}: List of estimated restitution coefficients for left-wall impacts.
  \item \code{restitution_right_estimates}: List of estimated restitution coefficients for right-wall impacts.
\end{itemize}

\paragraph{Environment-specific detectability rules}
\begin{itemize}
  \item If the class is \code{viscous_damping} or \code{mass} and the inclination is below 0.5, then there must be a right wall contact before and after \interventionTerm
  \item The total number of impacts on any wall must be less than 20.
\end{itemize}

{\large\bfseries Plot\par}

\begin{figure}[ht]
\centering
\IfFileExists{plots/environments/b054165364e7ab8827bbcaf351c56534__a711d5ad4b1191cbb8bf4e986e311740__Initial_Spacer_F_x__noise_stack_low_high.pdf}{%
  \includegraphics[width=0.85\linewidth]{plots/environments/b054165364e7ab8827bbcaf351c56534__a711d5ad4b1191cbb8bf4e986e311740__Initial_Spacer_F_x__noise_stack_low_high.pdf}%
}{%
  \includegraphics[width=0.85\linewidth]{docs/plots/environments/b054165364e7ab8827bbcaf351c56534__a711d5ad4b1191cbb8bf4e986e311740__Initial_Spacer_F_x__noise_stack_low_high.pdf}%
}
\caption{Position channel from DampedMassBetweenWalls.}
\label{fig:balldrop_env}
\end{figure}

\par\bigskip
{\LARGE\bfseries\texttt{InclinedPlane}\par}
\medskip
{\large\bfseries Textual-description levels\par}
\textbf{Low textual context}
\begin{DocCode}
These time series were generated by simulating a physical process.
\end{DocCode}

\textbf{High textual context}
\begin{DocCode}
These time series were generated by simulating a block of mass m moving along an infinitely long rigid plane inclined at an angle theta under gravitational acceleration g and Coulomb friction.
The coordinate axis is aligned with the plane. 
The block is also subject to an externally applied periodic force along the plane.
The friction force acts along the plane and opposes motion.
When the block is moving, that is, when v(t) is nonzero, friction is modeled as kinetic Coulomb friction.
A breakaway static-friction threshold is also modeled: when the block is at rest, motion starts only if the net driving force along the plane exceeds a breakaway limit.
\end{DocCode}

\begin{EnvironmentConfigTableNoMod}
\textemdash & \code{initial_velocity} & \code{[-5, 12]} \\
\textemdash & \code{Amplitude} & \code{[-200, 200]} \\
\textemdash & \code{Period} & \code{[0.5, 3]} \\
\hline
gravity acceleration & \code{gravity_constant} & \code{[2, 28]} \\
plane inclination angle & \code{plane_inclination} & \code{[5, 70]} \\
mass & \code{mass} & \code{[0.5, 10]} \\
Coulomb friction coefficient & \code{coulomb_friction_coefficient} & \code{[0.2, 3]} \\
breakaway friction coefficient & \code{breakaway_friction_coefficient} & \code{[0.2, 15]} \\
\end{EnvironmentConfigTableNoMod}

\begin{ParameterSrdTable}
\code{gravity_constant} & \code{0.3} & \code{0.06} \\
\code{plane_inclination} & \code{0.3} & \code{0.06} \\
\code{mass} & \code{0.3} & \code{0.06} \\
\code{coulomb_friction_coefficient} & \code{0.3} & \code{0.06} \\
\code{breakaway_friction_coefficient} & \code{0.3} & \code{0.06} \\
\end{ParameterSrdTable}

\begin{SimulationPropertiesTable}{Simulation Properties}
\code{end_time_input_s} & \code{10.0} \\
\code{sampling_rate_hz} & \code{50.0} \\
\end{SimulationPropertiesTable}

\begin{MappingTable}{Observed Signals.}
col1: velocity of the mass along the plane & \code{mass_velocity} \\
col2: friction force & \code{friction_force} \\
col3: normal force & \code{normal_force} \\
col4: time & \code{time} \\
\end{MappingTable}

\paragraph{Problem-specific features.}
\begin{itemize}
  \item \code{moving_intervals}
  \item \code{static_intervals}
  \item \code{static_to_moving_transition_times}
  \item \code{moving_to_static_transition_times}
  \item \code{normal_force_pre_post_shift}
  \item \code{kinetic_friction_ratio_pre_post_shift}
  \item \code{post_intervention_acceleration_shift}
  \item \code{periodic_force_response_amplitude}
\end{itemize}

\paragraph{Environment-specific detectability rules}
\begin{itemize}
  \item If the class is \code{breakaway_friction_coefficient} then \termMinSRDDistance must be higher than 0.2 for at least 5 samples in a row.  
  \item If the correct class is \code{gravity_constant} or \code{mass} then velocity must be different than 0 before and after the \interventionTerm for at least 5 samples before and 5 samples after
\end{itemize}

\newpage
{\large\bfseries Plot\par}

\begin{figure}[ht]
\centering
\IfFileExists{plots/environments/02fccb048ef7fb32d53361dc56ac8748__coulomb_friction_coefficient__mass_velocity__noise_stack_low_high.pdf}{%
  \includegraphics[width=0.85\linewidth]{plots/environments/02fccb048ef7fb32d53361dc56ac8748__coulomb_friction_coefficient__mass_velocity__noise_stack_low_high.pdf}%
}{%
  \includegraphics[width=0.85\linewidth]{docs/plots/environments/02fccb048ef7fb32d53361dc56ac8748__coulomb_friction_coefficient__mass_velocity__noise_stack_low_high.pdf}%
}
\caption{Velocity channel from InclinedPlane.}
\label{fig:balldrop_env}
\end{figure}
\FloatBarrier

\tsenvAppendixEnvironmentEnd
